\textbf{Datos y mercado.} Se asume que los precios de cierre y dividendos de los CSV representan adecuadamente la informaci\'on disponible para cada activo, por lo que el retorno total incorpora reinversi\'on impl\'icita de dividendos. La matriz de retornos se construye con \texttt{join="inner"}, suponiendo que eliminar fechas no comunes no introduce un sesgo relevante. La metadata de capitalizaci\'on de mercado se considera suficiente para construir el benchmark y los pesos de equilibrio de Black-Litterman, aunque puede existir sesgo de supervivencia. Las carteras son \textit{long-only}, cumplen $w_i \geq 0$ y $\sum w_i = 1$, no usan apalancamiento y permiten divisibilidad perfecta de acciones. La optimizaci\'on omite costos de transacci\'on, impuestos, \textit{spreads} bid-ask, restricciones de liquidez, impacto de mercado y \textit{slippage}. La utilidad de FinPUC se modela como una comisi\'on mensual de 0,5\% sobre el saldo activo administrado, sin comisiones por transacci\'on ni por rebalanceo. La tasa libre de riesgo se mantiene constante en $R_f = 2\%$ anual.

\textbf{Supuestos estad\'isticos.} Los retornos hist\'oricos se usan como informaci\'on para estimar retornos y riesgos futuros, bajo estabilidad temporal parcial de $\mu$ y $\Sigma$. Para estabilizar la covarianza muestral se aplica \textit{shrinkage} de 20\%: $\Sigma_{\text{shrink}} = (1-\lambda)\Sigma + \lambda \cdot \text{diag}(\Sigma)$, con $\lambda = 0.20$. La anualizaci\'on considera 252 d\'ias h\'abiles y escala la volatilidad por $\sqrt{252}$. En la simulaci\'on Monte Carlo, los retornos semanales se aproximan mediante $\mathcal{N}(\mu_{\text{sem}},\, \sigma_{\text{sem}}^2)$, aun cuando el an\'alisis distribucional muestra colas pesadas y curtosis elevada en los retornos diarios.

\textbf{Optimizaci\'on y rebalanceo.} Los problemas formulados en CVXPY se tratan como convexos y adecuados para representar la decisi\'on de inversi\'on. Si una cota de volatilidad vuelve infactible un perfil, se relaja solo esa restricci\'on. El portafolio de tangencia se aproxima eligiendo el mayor Sharpe dentro de la frontera eficiente discretizada, usando 30 puntos en Markowitz y 12 en Black-Litterman. En Black-Litterman, las \textit{views} no requieren cubrir todos los activos: aquellos sin \textit{view} expl\'icita conservan el prior de equilibrio $\pi$ y reciben ajustes indirectos por correlaci\'on a trav\'es de $\Sigma$. Las variantes por capitalizaci\'on usan subconjuntos de 20 a 40 activos, mientras que momentum general y desempleo consideran los 601 activos. El rebalanceo se realiza cada 126 d\'ias h\'abiles; entre rebalanceos los pesos permanecen constantes y, en cada fecha, la ventana de entrenamiento incorpora solo retornos ya observados.

\textbf{Experimento pandemia.} La comparaci\'on entre escenarios ``con pandemia'' y ``sin pandemia'' mantiene fijo el per\'iodo futuro fuera de muestra. Por ello, las diferencias se interpretan como efecto de la informaci\'on utilizada para calibrar los modelos, no como cambios en el mercado futuro. La base sin pandemia excluye las observaciones de 2020--2022, mientras que la base con pandemia las conserva.

\textbf{Simulaci\'on Monte Carlo (P4).} Cada cliente inicia con $C_0 = 1\,000$ USD. La utilidad de FinPUC proviene exclusivamente de una comisi\'on mensual de 0,5\% sobre el saldo activo administrado, acumulada en la implementaci\'on semanal con tasa equivalente $k_{\text{sem}} = 0{,}005 \times 12/52$; no hay comisiones por transacci\'on ni por rebalanceo. El comportamiento del cliente se reduce a una probabilidad de abandono semanal (modelo \texttt{views\_v1\_plus}): la p\'erdida se mide respecto de $C_0$ y la probabilidad conductual sigue una log\'istica sin reescalar, $1/(1+\exp(-(x_t-\hat{x})))$, activa solo cuando la p\'erdida supera la tolerancia $\hat{x}$ del perfil; adem\'as, si la riqueza llega a cero, el abandono es forzoso (restricci\'on de ruina). Para el perfil muy conservador, con tolerancia 0\%, cualquier p\'erdida genera probabilidad de abandono positiva. No se modela una probabilidad separada de aceptaci\'on de recomendaciones; la activaci\'on inicial del cliente se trata como mec\'anica de entrada. Se ejecutan 2\,000 trayectorias por combinaci\'on, cantidad considerada adecuada para estabilizar riqueza terminal, tasa de abandono y utilidad de la empresa.

\textbf{Sensibilidad.} El impacto de los supuestos se eval\'ua mediante tres ejercicios: comparaci\'on del entrenamiento con y sin pandemia, midiendo el delta de Sharpe por portafolio y t\'ecnica; validaci\'on con 70 splits aleatorios y 18 ventanas m\'oviles robustas, con horizontes desde h2\_f7 hasta h7\_f2; y proyecciones Monte Carlo desfavorable, neutra y favorable, modificando el retorno esperado en $\pm 0.5\sigma$ para observar cambios en riqueza terminal, tasa de retiro y utilidad de la empresa.

\textbf{Supuestos de la Entrega 3.}
La Tabla~\ref{tab:supuestos_e3} resume los supuestos espec\'ificos de la tercera fase del proyecto, adicionales a los ya establecidos en la Entrega 2.

\begin{table*}[h]
\centering
\caption{Supuestos adicionales de la Entrega 3.}
\label{tab:supuestos_e3}
\scriptsize
\begin{tabular}{@{}ll@{}}
\toprule
Supuesto & Valor \\
\midrule
Caso base & Markowitz (no equiponderado) \\
Objetivo de recomendaci\'on & Maximizar el retorno esperado del cliente \\
Comisi\'on de FinPUC & 0,5\% mensual sobre el saldo activo administrado \\
Tasa semanal equivalente & $k_{\text{sem}} = 0{,}005 \times 12/52$ \\
Frecuencia de recomendaci\'on / abandono & Semestral / semanal \\
Modelo de abandono & \'Unico, \texttt{views\_v1\_plus} (log\'istica de umbral + ruina) \\
Probabilidad de abandono & $1/(1+\exp(-(x_t-\hat{x})))$ si $x_t>\hat{x}$; $0$ si no \\
Restricci\'on de ruina & Abandono forzoso si riqueza $\leq 0$ \\
Aceptaci\'on de recomendaci\'on & No se modela; activaci\'on inicial como mec\'anica de entrada \\
Calibraci\'on BL & Independiente por \textit{view} (config propia por metodolog\'ia) \\
Retornos semanales en Monte Carlo & Normales i.i.d.\ con KPIs calibrados por metodolog\'ia \\
Trayectorias por combinaci\'on & 2\,000 \\
Horizontes de validaci\'on & $\mathcal{H}_1$ calibraci\'on, $\mathcal{H}_2$ validaci\'on, $\mathcal{H}_3$ test limpio P4 \\
\bottomrule
\end{tabular}
\end{table*}
